\section{到 Borel 的有理约化空间}
\markboth{到 Borel 的有理约化空间}{}

众所周知，叠 $\Bun_B$ 不连通。因此，若想用伴随函子对
\[
  \fP_! : D(\Bun_B)\rightleftarrows D(\Bun_G) : \fP^!
\]
来逼近范畴 $D(\Bun_G)$，就会丢失大量信息。

为此我们转向 Drinfeld 紧化 $\bBun_B$ 及其映射
\[
  \bfP : \bBun_B \to \Bun_G .
\]
然而 $\bBun_B$ 另有一个毛病：冗余太多。也就是说，它不只是把 $\Bun_B$ 的各个连通分支彼此粘起来，而是沿着曲线的部分对称化幂把各分支的乘积粘起来。

本节的目标是定义一个「空间」$\Bun\rt_B$，它把 $\Bun_B$ 的所有连通分支真正无冗余地粘在一起。唯一的问题是：这个空间不存在。不过范畴 $D(\Bun\rt_B)$ 存在，而这正是我们要做的事。想法很简单：考察 $D(\bBun_B)$，再用关于适当群胚的等变性把冗余消去。

\ob{1.1}{blk:1.1}\textbf{「空间」。}考虑叠 $\bBun_B$。令 $\mathcal H_{G,B}$ 为群胚
\[
  \mathcal H_{G,B}\subset \bBun_B\underset{\Bun_G}{\times}\bBun_B ,
\]
它由给定 $G$-丛上、在曲线的一般点处相等的广义 $B$-约化对构成。

我们把范畴 $D(\Bun\rt_B)$ 定义为 $D(\bBun_B)$ 关于 $\mathcal H_{G,B}$ 的等变范畴。

\emph{注。}可以证明 $\Bun\rt_B$ 不作为 Artin 叠存在，即不存在一个叠，带有映射 $\Bun_B\to\Bun\rt_B$ 且该映射诱导出遗忘函子 $D(\Bun\rt_B)\to D(\Bun_B)$。\footnote{不存在的证明来自：$D(\Bun\rt_B)$ 上不存在具有预期性质的 $t$-结构。}然而为记号方便，我们仍把它当作这样的叠来处理。

\begin{annote}
「广义 $B$-约化」是 Drinfeld 紧化 $\bBun_B$ 的原始定义中的对象：在每个点处允许 $B$-约化带有极点，极点由一个 $G$-支配余特征标刻画。两点在一般点处相等，等价于说两者的差别由曲线的除子给出；这些差别正是下面 $X^\lambda$ 所参数化的冗余。\par\vspace{0.15em}{\raggedleft 【来源：原书 p.2 §1.1（A）；[GL:contr] 的引言（A）】\par}
\end{annote}

\ob{1.2}{blk:1.2}\textbf{单子描述。}令 $\oblv_B$ 记遗忘函子
\[
  D(\Bun\rt_B)\to D(\bBun_B).
\]
投影
\[
  \bBun_B\xleftarrow{h^\leftarrow}\mathcal H_{G,B}\xrightarrow{h^\rightarrow}\bBun_B
\]
都是 ind-紧合的。因此函子 $\oblv_B$ 有左伴随，记作 $\ind_B$。复合
\[
  \oblv_B\circ\ind_B : D(\bBun_B)\to D(\bBun_B)
\]
等同于 $\Av_{\mathcal H_{G,B}} := h^\rightarrow_*\circ h^!_\leftarrow$。由 Barr-Beck，我们可以把 $D(\Bun\rt_B)$ 看作 $D(\bBun_B)$ 上关于单子 $\Av_{\mathcal H_{G,B}}$ 的模。

\begin{annote}
为何 ind-紧合给出左伴随：对 ind-紧合映射，$!$-拉回与 $*$-推前构成伴随对；等变范畴的遗忘函子是取关于该群胚的「平均」，而平均由 $h^\rightarrow_*\circ h^!_\leftarrow$ 实现，故遗忘函子的左伴随是沿 $h^\rightarrow$ 的 $!$-推前。Barr-Beck 定理则把等变范畴等同于该单子的模范畴。\par\vspace{0.15em}{\raggedleft 【来源：原书 p.2 §1.2（A）；[GL:DG] 关于 D-模六函子形式的约定（A）】\par}
\end{annote}

\ob{1.3}{blk:1.3}\textbf{与 $\Bun_G$ 的关系。}有伴随函子对
\[
  \bfP_! : D(\bBun_B)\rightleftarrows D(\Bun_G) : \bfP^! .
\]
函子 $\bfP^!$ 典范地分解为 $\oblv_B\circ(\fP\rt)^!$，其中 $(\fP\rt)^!$ 是典范定义的函子 $D(\Bun_G)\to D(\Bun\rt_B)$。于是 $(\fP\rt)^!$ 有左伴随，记作 $(\fP\rt)_!$，使得
\[
  \bfP \simeq (\fP\rt)_!\circ\ind_B .
\]
下述将在 [GL:contr] 中证明：

\begin{othm}{拟定理}{1.3.1}{thm:1.3.1}
函子 $(\fP\rt)^!$ 全忠实。
\end{othm}

\begin{annote}
分解 $\bfP^!\simeq\oblv_B\circ(\fP\rt)^!$ 的含义：先在 $\Bun_G$ 上取 $!$-拉回，再沿 $\Bun\rt_B\to\Bun_G$ 的「部分定义」拉回；后者之所以存在，是因为 $\bfP^!$ 的像落在 $\mathcal H_{G,B}$-等变对象里。由此 $(\fP\rt)^!$ 是 $\bfP^!$ 在等变范畴中的「提升」。拟定理 1.3.1 说这一提升不丢失信息——它是本文反复调用的基本事实（§7.5 的 Quasi-Theorem 7.5.3 是其抛物推广）。\par\vspace{0.15em}{\raggedleft 【来源：原书 p.2 §1.3（A）；译者按（上下文推断，无外部依据）】\par}
\end{annote}
